\section{Threats to Validity and Limitations}\label{sec:threats}

\subsection{Construct Validity}\label{sec:threats_construct}
\textbf{Calendar duration vs.\ debugging effort.} The central construct threat is the distinction between observed calendar duration and developer debugging effort~\cite{mockus2002two}. The duration $\texttt{closed\_at} - \texttt{created\_at}$ measures elapsed wall-clock calendar time. In open-source software, an issue that remains open for 100~days may involve only two hours of actual developer investigation separated by 99~days of inactivity, triage queueing, or waiting for feedback. Conversely, an issue closed in 48~hours may have occupied three core maintainers in intensive debugging. We explicitly state: \emph{no finding in this report measures developer effort, work hours, or cognitive difficulty}.

\textbf{Label validity.} The cohort is defined by maintainer labels, and the validation audit (\S\ref{sec:validation}) estimates only \emph{precision}: what fraction of what we called a bug really is one. We have no estimate of \emph{recall}---how many genuine bugs carry no strict label---and the audit was an assisted single-reviewer screening, not two independent human reviewers. Projects that never use a \texttt{bug}-type label are under-represented in the cohort, and the sample of $\valvN$ issues yields confidence intervals several points wide.

\textbf{Issue closure vs.\ successful defect resolution.} A fundamental assumption in repository studies is that closing an issue signifies that the defect was successfully resolved. In reality, bugs are frequently closed without a fix: issues may be closed due to obsolescence, inability to reproduce, duplicate reporting, or project abandonment~\cite{herzig2013not}. Although our cohort criteria require at least one linked pull request—and robustness check R3 confirms that restricting to issues with at least one \emph{merged} PR yields identical conclusions ($\delta = 0.283$, 96.5\% of Group~B merged)—a merged PR does not guarantee that the defect was completely eliminated without introducing regressions.

\textbf{PR activity metrics as proxies.} In exploratory analyses (\S\ref{sec:exploratory}), we examine commit counts, comments, and contributor numbers. These metrics represent repository interaction activity rather than direct measurements of bug complexity. High comment counts may reflect contentious code reviews or community chatter rather than algorithmic difficulty.

\subsection{Internal Validity and Exposure Bias}\label{sec:threats_internal}
\textbf{Time-dependent link accumulation and landmark limitations.} The primary methodological threat is reverse causality via link accumulation: open issues remain available to receive pull requests while closed issues cannot, creating a positive dependency between lifespan and PR count. Our timing analysis (\S\ref{sec:timing}) confirms that the last linked PR in Group~B arrives at a median of \valtMedLast\% of total lifetime. The fixed-landmark analysis (\S\ref{sec:landmark}) reduces this calendar accumulation bias by fixing group membership at day~30, where the observed difference is $-0.8$ days ($p=0.832$). However, several critical caveats apply:
\begin{itemize}
  \item \emph{Estimation imprecision}: With only 33 multi-PR issues surviving to day~30, the test has low power and the wide 95\% CI ($-22.5$ to $+29.2$ days) cannot rule out meaningful differences in either direction. The result is inconclusive rather than proof of absence.
  \item \emph{Survivor and cohort selection}: The landmark analysis is restricted to issues that eventually closed and eventually had at least one linked PR; tickets still open at the snapshot or never linked to a PR are excluded. Landmarking mitigates time-dependent exposure bias within closed issues, but does not eliminate all selection bias.
  \item \emph{Landmark sensitivity}: Results differ across milestones: at day~7, multi-PR issues close faster ($-5.1$ days, $p=0.0154$), likely capturing early concerted maintainer focus, while at day~90 the difference is $+21.2$ days but statistically non-significant ($p=0.641, n=22$).
\end{itemize}
We caution that neither the raw $\valmedDiff$-day median gap nor the landmark estimates represent definitive causal effects.

\textbf{Unrecorded link timestamps.} As flagged in \S\ref{sec:inventory} and \S\ref{sec:landmarkdesign}, the \texttt{issue\_pr} table stores static issue--PR links without cross-reference event timestamps. Consequently, PR creation time must serve as a proxy for link timing: a PR created early but linked late still counts at the landmark, and no exact prospective landmark model is possible with this data.

\textbf{Snapshot censoring.} Only issues with a recorded \texttt{closed\_at} enter the analysis, so bugs still open at the snapshot date (\valsnapshot) are absent—plausibly the slowest, hardest ones. The exploratory survival analysis in \S\ref{sec:exploratory} brackets this, but its link counts for open issues are truncated by the snapshot and can only grow, so it cannot fully repair the bias. All reported durations are therefore lower bounds on true resolution times.

\subsection{External Validity}\label{sec:threats_external}
\textbf{Domain specificity of game software.} OSSGameBench samples exclusively from open-source video games and game development engines~\cite{marsad2026ossgamebench}. Game development differs from traditional enterprise software due to real-time performance constraints, tight coupling between compiled engines and binary assets (textures, meshes, shaders, audio), and reliance on volunteer game modders. Findings may not generalize to enterprise systems, microservices, or developer tooling.

\textbf{Platform and governance bias.} The corpus derives exclusively from GitHub repositories. Open-source projects operating on self-hosted GitLab, Gerrit, or Bugzilla instances with different code review conventions are not represented.
